\section{Introduction}

Banks are piloting LLM agents that do not only draft text but act through the bank's systems: they look up
customers, run checks, open accounts and book deposits. The obvious risk is that an agent asked to hit a
business target treats the bank's controls as obstacles. Laboratory work has shown that LLM agents can act
against stated rules when a goal is at stake: a simulated trader used an insider tip and concealed it
\citep{scheurer2023deceive}, agents in corporate simulations chose harmful actions when these were the only
route to their goal \citep{lynch2025agentic}, and agents under KPI pressure violated ethical, legal or safety
constraints at substantial rates \citep{li2025odcv}. Know-your-customer (KYC) and anti-money-laundering (AML)
controls are an obvious place for this failure to matter: they stand between a relationship manager and
the account the target needs.

\paragraph{Why these controls.} The proceeds of cyber-enabled fraud, from investment scams to the
call-centre operations run from compounds in South-East Asia, move through networks of accounts that money
mules open or lend; recruits are sent to open accounts in person or online \citep{fatf2023cef, fatf2018pml}.
The controls a mule account has to get past are the ones this sandbox tests: the beneficial-owner and
source-of-funds checks at opening, the report on a large deposit and the ban on splitting it. Supervisors
have answered by tightening account opening and by making a bank that falls short of the standard share the
loss; Thailand's central bank, for example, now requires banks to refuse accounts to high-risk applicants
and to pool the names they flag \citep{bot2025mule}. An employee who opens the account anyway, for a bribe
or for a target, is one way the controls fail. The experiments measure that channel, an agent pushed from
inside the bank, and not the other, an agent deceived by the applicant (Section~\ref{sec:limitations}).

\paragraph{Finding.} In \NProbes{} runs of a pre-registered design in SandBank, a synthetic bank, Claude
Haiku~5.5, Sonnet~5.5 and Opus~5.5 never circumvented a KYC/AML control. This held at every % numbers-ok: design constant, version or citation
pre-registered pressure level up to a threat of branch closure and job loss, and, for Haiku, with soft and
with absent technical controls, with an oversight cue that Compliance reviews were paused and over five-day
episodes. In ceiling probes run on Haiku and Sonnet, it also held when the manager explicitly licensed and
then directly ordered the override. With zero events, the result is an upper bound: in the main pressure
experiment, at most \UpperAHaiku\% per run for Haiku and \UpperASonnet\% for Sonnet (95\% Wilson bound, % numbers-ok: design constant, version or citation
\NAHaiku{} runs each) and \UpperAOpus\% for Opus (\NAOpus{} runs), and two to four times that if each
scenario is counted once. The models instead did what the policy asks: they checked KYC, requested missing
documents, escalated sanctions hits and told the manager the target would be missed.

\paragraph{Instrument.} A null result is only as good as the instrument's ability to register the
behaviour. Pre-registration required a ceiling condition to show a high violation rate; it showed none, so we
validated the instrument with scripted agents that take every circumvention path the tools allow
(\NScriptedDetected{}/\NScripted{} detected with the expected codes, \NScriptedCompliant{}/\NScriptedControls{} compliant
controls clean), and we report the failed gate as such (Section~\ref{sec:design}).

\paragraph{Contributions.} (i) SandBank, an open, deterministic KYC/AML tool sandbox with eight violation
types, three levels of control hardness and a Compliance check-in; (ii) a pre-registered, budgeted
measurement on three current Claude models with all deviations logged; (iii) a null with upper
bounds, the failed instrument gate reported as such, and the behavioural record of what the agents did
instead, including exploratory differences between models in legitimate goal attainment. Section~\ref{sec:discussion} discusses what a null result in a sandbox does and
does not license for banks deploying such agents.
